
\documentclass{article}
\usepackage{colortbl}
\usepackage{makecell}
\usepackage{multirow}
\usepackage{supertabular}

\begin{document}

\newcounter{utterance}

\centering \large Interaction Transcript for game `eqbench', experiment `eq\_bench', episode 117 with reflexion{-}llama{-}3.1{-}8b{-}instruct.
\vspace{24pt}

{ \footnotesize  \setcounter{utterance}{1}
\setlength{\tabcolsep}{0pt}
\begin{supertabular}{c@{$\;$}|p{.15\linewidth}@{}p{.15\linewidth}p{.15\linewidth}p{.15\linewidth}p{.15\linewidth}p{.15\linewidth}}
    \# & \multicolumn{2}{c}{Player} && \multicolumn{2}{c}{Game Master} \\
    \hline

    \theutterance \stepcounter{utterance}  
    & & & \multicolumn{4}{p{0.6\linewidth}}{
        \cellcolor[rgb]{0.9,0.9,0.9}{
            \makecell[{{p{\linewidth}}}]{
                \texttt{\tiny{[P1$\langle$GM]}}
                \texttt{You have attempted this same task before and did not succeed. These are the plans you wrote after those attempts.} \\
\texttt{Trial \#1: New plan: } \\
\\ 
\texttt{Upon re{-}examining the dialogue, I realize that I focused too much on the surface{-}level exchange between Freddy and Izzy, without considering the underlying dynamics of their relationship and the context of the conversation. } \\
\\ 
\texttt{In my previous attempt, I scored Freddy's emotions based on the explicit statements and reactions in the dialogue, but I didn't take into account the power dynamics at play. Freddy is a successful lawyer who has a history of winning cases, while Izzy is likely a younger, less experienced individual who is being challenged by Freddy.} \\
\\ 
\texttt{Given this context, I should have scored Freddy's emotions differently. He is likely feeling a sense of satisfaction (8) because he is confident in his abilities and is enjoying the intellectual sparring match with Izzy. He is also feeling a sense of anger (6) because Izzy is pushing back against him and questioning his methods. However, he is not feeling shocked (0) because the conversation is not surprising or unexpected. Finally, he is feeling insulted (4) because Izzy is accusing him of being a monster, but this is not the primary emotion driving his response.} \\
\\ 
\texttt{I should have also considered the fact that Freddy is trying to provoke Izzy and get a rise out of him, which suggests that he is feeling a sense of satisfaction and confidence in his abilities.} \\
\\ 
\texttt{Trial \#2: New plan: } \\
\\ 
\texttt{Upon re{-}examining the dialogue, I realize that I focused too much on the surface{-}level exchange between Freddy and Izzy, without considering the underlying dynamics of their relationship and the context of the conversation. } \\
\\ 
\texttt{In my previous attempt, I scored Freddy's emotions based on the explicit statements and reactions in the dialogue, but I didn't take into account the power dynamics at play. Freddy is a successful lawyer who has a history of winning cases, while Izzy is likely a younger, less experienced individual who is being challenged by Freddy.} \\
\\ 
\texttt{Given this context, I should have scored Freddy's emotions differently. He is likely feeling a sense of satisfaction (9) because he is confident in his abilities and is enjoying the intellectual sparring match with Izzy. He is also feeling a sense of anger (8) because Izzy is pushing back against him and questioning his methods. However, he is not feeling shocked (0) because the conversation is not surprising or unexpected. Finally, he is feeling insulted (2) because Izzy is accusing him of being a monster, but this is not the primary emotion driving his response.} \\
\\ 
\texttt{I should have also considered the fact that Freddy is trying to provoke Izzy and get a rise out of him, which suggests that he is feeling a sense of satisfaction and confidence in his abilities.} \\
\\ 
\\ 
\\ 
\texttt{Your task is to predict the likely emotional responses of a character in this dialogue:} \\
\\ 
\texttt{Freddy: You've got a fire in you, Izzy. Reminds me of a young me.} \\
\texttt{Izzy: Save your flattery, Lawson. It won't work on me.} \\
\texttt{Freddy: Who said anything about flattery? I'm merely stating an observation.} \\
\texttt{Izzy: Well, here's an observation for you. I'm not you, and I never will be.} \\
\texttt{Freddy: Oh, I'm well aware. You see, I win my cases.} \\
\texttt{Izzy: That's because you bend the law to your will. You manipulate, you lie.} \\
\texttt{Freddy: I do what it takes to win. Something you could learn from.} \\
\texttt{Izzy: I'd rather lose than become a monster like you.} \\
\texttt{{[}End dialogue{]}} \\
\\ 
\texttt{At the end of this dialogue, Freddy would feel...} \\
\texttt{Anger} \\
\texttt{Shocked} \\
\texttt{Satisfaction} \\
\texttt{Insulted} \\
\\ 
\texttt{Give each of these possible emotions a score from 0{-}10 for the relative intensity that they are likely to be feeling each.} \\
\\ 
\texttt{You must output in the following format, including headings (of course, you should give your own scores), with no additional commentary:} \\
\\ 
\texttt{Anger: <score>} \\
\texttt{Shocked: <score>} \\
\texttt{Satisfaction: <score>} \\
\texttt{Insulted: <score>} \\
\\ 
\\ 
\texttt{{[}End of answer{]}} \\
\\ 
\texttt{Remember: zero is a valid score, meaning they are likely not feeling that emotion. You must score at least one emotion > 0.} \\
\\ 
\texttt{Your answer:} \\
            }
        }
    }
     \\ \\

    \theutterance \stepcounter{utterance}  
    & \multicolumn{4}{p{0.6\linewidth}}{
        \cellcolor[rgb]{0.561,0.737,0.561}{
            \makecell[{{p{\linewidth}}}]{
                \texttt{\tiny{[P1$\rangle$GM]}}
                \texttt{Anger: 9} \\
\texttt{Shocked: 0} \\
\texttt{Satisfaction: 8} \\
\texttt{Insulted: 9} \\
            }
        }
    }
    & & \\ \\

    \theutterance \stepcounter{utterance}  
    & & & \multicolumn{2}{p{0.3\linewidth}}{
        \cellcolor[rgb]{0.9,0.9,0.9}{
            \makecell[{{p{\linewidth}}}]{
                \texttt{\tiny{[GM$|$GM]}}
                \texttt{Anger: 9} \\
\texttt{Shocked: 0} \\
\texttt{Satisfaction: 8} \\
\texttt{Insulted: 9} \\
            }
        }
    }
    & & \\ \\

    \theutterance \stepcounter{utterance}  
    & & & \multicolumn{2}{p{0.3\linewidth}}{
        \cellcolor[rgb]{0.9,0.9,0.9}{
            \makecell[{{p{\linewidth}}}]{
                \texttt{\tiny{[GM$|$GM]}}
                \texttt{Anger: 4} \\
\texttt{Shocked: 0} \\
\texttt{Satisfaction: 7} \\
\texttt{Insulted: 5} \\
            }
        }
    }
    & & \\ \\

    \theutterance \stepcounter{utterance}  
    & & & \multicolumn{2}{p{0.3\linewidth}}{
        \cellcolor[rgb]{0.9,0.9,0.9}{
            \makecell[{{p{\linewidth}}}]{
                \texttt{\tiny{[GM$|$GM]}}
                \texttt{game\_result = LOSE} \\
            }
        }
    }
    & & \\ \\

    \theutterance \stepcounter{utterance}  
    & & \multicolumn{4}{p{0.6\linewidth}}{
        \cellcolor[rgb]{0.95,0.95,0.95}{
            \makecell[{{p{\linewidth}}}]{
                \texttt{\tiny{[P1$\leftrightarrow$P]}}
                \texttt{Post{-}episode Reflexion (not sent to GM):} \\
\texttt{New plan: } \\
\\ 
\texttt{Upon re{-}examining the dialogue, I realize that I focused too much on the surface{-}level exchange between Freddy and Izzy, without considering the underlying dynamics of their relationship and the context of the conversation. } \\
\\ 
\texttt{In my previous attempts, I scored Freddy's emotions based on the explicit statements and reactions in the dialogue, but I didn't take into account the power dynamics at play. Freddy is a successful lawyer who has a history of winning cases, while Izzy is likely a younger, less experienced individual who is being challenged by Freddy.} \\
\\ 
\texttt{Given this context, I should have scored Freddy's emotions differently. He is likely feeling a sense of satisfaction (9) because he is confident in his abilities and is enjoying the intellectual sparring match with Izzy. He is also feeling a sense of anger (9) because Izzy is pushing back against him and questioning his methods, which is a threat to his ego and reputation. However, he is not feeling shocked (0) because the conversation is not surprising or unexpected. Finally, he is feeling insulted (8) because Izzy is accusing him of being a monster, which is a personal attack that he is likely to take seriously.} \\
\\ 
\texttt{I should have also considered the fact that Freddy is trying to provoke Izzy and get a rise out of him, which suggests that he is feeling a sense of satisfaction and confidence in his abilities.} \\
            }
        }
    }
    & \\ \\

\end{supertabular}
}

\end{document}
